\documentclass[12pt]{article}
\usepackage[utf8]{inputenc}
\usepackage[T1]{fontenc}
\usepackage{hyperref}
\title{Zhangzhung Language Resources for K-12 Curriculum Development}
\author{Compiled by Claude Sonnet 5.5 (claude-sonnet-5-5) for Andr\'as Kornai}
\date{}
\begin{document}
\maketitle

\section{Introduction}
Zhangzhung (ISO 639-3 xzh) was a Tibeto-Burman language of western Tibet. It is extinct, and it survives in Bon religious texts and in a few manuscripts. The search results contained scholarly materials only. No school materials, no community revival programme and no teaching resources were found.

\section{Primary Educational Resources}
None were found.

\section{Linguistic and Reference Materials}
\begin{itemize}
\item Erik Haarh, \emph{The Zhang-zhung Language: A Grammar and Dictionary of the Unexplored Language of the Tibetan Bonpos} (Acta Jutlandica, 1968). Library catalogue records:
\url{https://www.oeaw.ac.at/resources/Record/990001864100504498?lng=en}, \url{https://www.appliedbuddhism.org.uk/node/29671}, \url{https://find.library.upenn.edu/catalog/998390953503681}
\item Reference record: \url{https://wals.info/refdb/record/Haarh-1968}
\item ``The Present Stage of Deciphering Old Zhangzhung'' (academic paper): \url{https://minpaku.repo.nii.ac.jp/record/2570/files/SES75_010.pdf}
\item Another academic paper from the same repository: \url{https://minpaku.repo.nii.ac.jp/record/7403/files/SER76_002.pdf}
\item Glottolog reference entry: \url{https://glottolog.org/resource/reference/id/302963}
\item Endangered Languages Project entry: \url{https://endangeredlanguages.com/node/33490}
\end{itemize}

\section{Religious and Community Resources}
The language was preserved by Bonpo priests as a sacred written language. No specific community resources were identified in the search results.

\section{Digital and Audio Resources}
None were found.

\section{Curriculum Development Context}
The language is known only from a limited corpus and is partly undeciphered, so spoken-language teaching is not realistic. A K-12 unit could instead be about Zhangzhung as a historical language, using Haarh's glossary and the Dunhuang-based scholarship. It would be best developed with Bon scholars.
\end{document}